En esta práctica se estudió cómo se dobla un rayo de luz al cruzar la frontera entre el aire y una lente de acrílico en forma de media luna. Con una fuente de luz, una rendija y una mesa de rayos graduada, se midió el ángulo del rayo transmitido para distintos ángulos de incidencia, primero haciendo entrar la luz por la cara plana de la lente y después por la cara curva. En el segundo caso la luz sale del acrílico hacia el aire, y a partir de los 45 grados dejó de salir por completo para reflejarse dentro de la lente. De la gráfica del seno de un ángulo contra el seno del otro se obtuvo un índice de refracción de 1.494 con la cara curva, prácticamente igual al valor de referencia de 1.49, y de 1.351 con la cara plana, 9.3 por ciento abajo. El ángulo crítico calculado fue de 42.0 grados. Coincide con lo observado.
